%% ===========================================================================
\section{Introduction}
\label{sec:intro}
%% ===========================================================================

Large language models (LLMs) are now routinely asked to act as \emph{judges}: given a context and a short list of options, they return the option they find best. This pattern underlies LLM-as-a-judge evaluation \citep{zheng2023mtbench,liu2023geval}, tool and function selection \citep{li2024fnctod}, and many agentic pipelines in which a model must decide whether to act, to do nothing, or to ask for clarification. Two properties matter as much as accuracy in such settings. The first is \emph{cost}: a judge sits inside an interactive loop and is called many times per user turn. The second is \emph{trustworthy confidence}: a system can act autonomously only if the judge's probability tells it reliably when it is likely to be wrong, so that uncertain cases can be deferred to the user or to a more expensive component \citep{geifman2017selective}. LLM judges are known to be sensitive to the order and labels of options \citep{wang2024fair,li2025calibraeval}, and their confidence must often be recalibrated before it can be used \citep{radharapu2026calibrating}.

This paper asks whether a small decision model, fine-tuned for the judging role and calibrated, can replace an open-weight LLM judge in such a loop. We study the question on a concrete and well-instrumented task: \emph{slot filling} in task-oriented dialogue. For a request such as ``a flight from \emph{Lyon} to \emph{Paris} on \emph{Thursday morning}'', the system must fill each slot of the current intent (departure city, arrival city, date, time\ldots) with the value expressed by the user, or leave it empty. Recent work relies on generative LLMs for this task \citep{hu2022icdst,heck2023chatgpt,li2024fnctod}. Generation is flexible, but it is costly, it can produce values that do not occur in the input, and the resulting scores are rarely usable to decide whether to act.

We instead separate two roles, following a long line of picklist-based dialogue state trackers \citep{lee2019sumbt,zhang2020dsdst} and of retrieve-then-rerank pipelines \citep{wu2020blink}:
\begin{enumerate}
  \item \textbf{Propose.} Inexpensive generators (a lexicon of training values, a zero-shot span extractor such as GLiNER \citep{zaratiana2024gliner}, the ontology, the dialogue history) build, for each slot, a short list of candidate values from the current context.
  \item \textbf{Judge.} A light judge either selects one candidate or answers ``not specified''. Its probability, once calibrated, drives a three-zone decision: \emph{act} (fill the slot), \emph{ignore} (leave it empty), or \emph{ask} (defer).
\end{enumerate}
Because the judge can only return one of the proposed candidates or abstain, it cannot hallucinate a value. Its errors therefore decompose cleanly into \emph{candidate recall} (the right value was never proposed), \emph{choice} (it was proposed but not selected) and \emph{abstention} (a value was selected for a slot the user did not express), which makes the architecture convenient to diagnose.

As the judge we use Laya \citep{laya2026}, a publicly released, non-autoregressive ``System~1'' decision model built on modern bidirectional encoders \citep{warner2025modernbert,marone2026mmbert}, which answers multiple-choice questions about a structured state in a single forward pass. We compare it, zero-shot and after fine-tuning, with three recent open-weight LLM judges of 27 to 180 billion parameters from two families \citep{qwen38,qwen38flashnext,gemma4report}, all prompted with the same options and scored from the log-probabilities of their first answer token. The evaluation spans three public benchmarks with complementary properties: MultiATIS++ \citep{xu2020multiatis} (one domain, single turn), MTOP \citep{li2021mtop} (eleven domains and many intents, single turn) and Multi3WOZ \citep{hu2023multi3woz} (multi-domain, multi-turn dialogue), in English and in French, with both oracle candidate lists (to measure the judge alone) and real ones (to measure the whole pipeline).

The general-AI question behind this setting is the trade-off between a large general model used zero-shot and a small model specialised for one decision. Recent studies show that small fine-tuned models can match or beat LLMs on classification and slot filling once a modest amount of labelled data is available \citep{pecher2025comparing,aguirre2025finetuning,rana2025zeroshot}. We add two dimensions that matter for deployment and that these studies do not measure: whether the judge \emph{knows when it does not know} (risk--coverage behaviour after calibration), and how robust it is when the shape of its input changes (list length, missing gold value, language).

\paragraph{Research questions} We organise the study around five questions:
\textbf{RQ1} (efficiency) Does a fine-tuned Laya judge approach an open-weight LLM judge in slot F1, at a much lower cost?
\textbf{RQ2} (decomposition) Which share of the error comes from candidate recall, and which from the choice, and how do they vary with the number of candidates $k$ and the rate of missing gold values?
\textbf{RQ3} (generalisation) Does the judge transfer from English to French?
\textbf{RQ4} (abstention) Do calibration and the decision rule yield a useful risk--coverage trade-off?
\textbf{RQ5} (training) Does random augmentation of the option lists at training time make the judge robust to a changing context?

\paragraph{Contributions}
\begin{enumerate}
  \item A \textbf{proposer--judge formulation} of slot filling as calibrated candidate selection with explicit abstention, and a reproducible conversion of three span-annotated benchmarks (two languages, single- and multi-turn) into this format, including a domain filter for multi-turn dialogue (\cref{sec:method,sec:setup}).
  \item A \textbf{fine-tuning recipe} for a non-autoregressive decision model: supervised optimisation of the log score on training items produced by the inference code itself, so that training and evaluation cannot differ in format, with checkpoint selection on development F1 and an \emph{option-list augmentation} designed to reproduce the confusions of real candidate generators (\cref{sec:method}).
  \item A \textbf{controlled comparison with three open-weight LLM judges} (27--180\,B parameters, two families) on identical inputs: the fine-tuned judge is more accurate on every benchmark, language and candidate type, 3 to 6 times faster per request and 5 to 38 times faster in batch throughput, mainly because it abstains correctly (\cref{sec:rq1}).
  \item A \textbf{reliability analysis} showing that after calibration the small judge's confidences rank its errors far better than those of the LLMs (area under the risk--coverage curve 3 to 40 times lower), so that it can decide most requests without asking at a few percent risk, which none of the LLM judges can do on MTOP (\cref{sec:rq4}).
  \item A \textbf{comparison with a fine-tuned sequence tagger} that delimits where the proposer--judge architecture pays off: the tagger alone is better when slot values are literal spans of the utterance, the pipeline when values come from the dialogue history or must be mapped to canonical forms; the judge itself is not the bottleneck, the proposers are (\cref{sec:rq2-tagger}).
  \item An \textbf{error decomposition and robustness study} across candidate recall, choice and abstention, list size, missing values and cross-lingual transfer, with a negative-then-corrected result on augmentation that shows \emph{which} augmentation helps (\cref{sec:rq2,sec:rq3,sec:rq5}), all under a reproducibility protocol with multiple seeds, cluster bootstrap intervals, Holm-corrected paired permutation tests and code fingerprints in every run (\cref{sec:stats}).
\end{enumerate}

The rest of the paper is organised as follows. \Cref{sec:related} reviews related work. \Cref{sec:method} formalises the proposer--judge architecture and describes the judges, the candidate generators, fine-tuning, augmentation and the calibrated decision. \Cref{sec:setup} details datasets, conditions, metrics and the statistical protocol. \Cref{sec:results} answers RQ1--RQ5, \cref{sec:discussion} discusses errors, cost, calibration, limitations and threats to validity, and \cref{sec:conclusion} concludes.

